% !TeX root = ../slides.tex
% 02-architectures.tex

\begin{frame}{Network architectures: state space encoding}
  \begin{itemize}
    \item Our original idea on how to represent the state space is based on the paper by
    \textcite{wang2025minesweepercnn}. The author used a \alert{one-hot board encoding} (10 channels: values $0$–$8$
    plus one channel for unrevealed cells) fed to a convolutional
    $Q$-network.
  \end{itemize}

  \resizebox{0.98\textwidth}{!}{
    \begin{tikzpicture}[font=\small]

      % Input tensor
      \node (input) at (0,0) {$[B,C,H,W]$};

      % CNN block
      \node[
        draw,
        rectangle,
        minimum width=4.2cm,
        minimum height=2.5cm,
        align=center
      ] (cnn) at (4.2,0) {
        {\Large CNN}\\[0.25cm]
        fully convolutional\\
        (or with global features)
      };

      % Output values
      \node[anchor=west] (outputs) at (8.4,0) {
        $\left\{
        \begin{aligned}
          Q(s,a_0)
          &\;/\;
          \ell(s,a_0)
          \\[0.25cm]
          &\vdots
          \\[0.25cm]
          Q(s,a_{H\cdot W-1})
          &\;/\;
          \ell(s,a_{H\cdot W-1})
        \end{aligned}
        \right.$
      };

      % Arrows
      \draw[->, thick]
        (input.east) -- (cnn.west);

      \draw[->, thick]
        (cnn.east)
        -- node[above] {$[B,H\cdot W]$}
        (outputs.west);

    \end{tikzpicture}
  }

  \begin{itemize}
    \item In a second stage of development we added a constant channel, the 10th one, containing
    the \textit{mine density}. This was done with the idea to let the network generalize well
    to harder games. In this case, $C = 11$.
  \end{itemize}
\end{frame}

\begin{frame}{Design Choices for Stable Training}
  \begin{itemize}
    \item \texttt{player\_board} content is not re-usable as PyTorch index; otherwise, we might have the risk to override the wrong one-hot channel.
    \begin{table}[h]
      \centering
      \begin{tabular}{cl}
        \textbf{One-hot channel ID} & \textbf{Encoded value / meaning} \\
        \midrule
        $0,\ldots,8$ & Values $0,\ldots,8$ of \texttt{player\_board} \\
        $9$          & $-2$: unrevealed cell \\
        $10$         & Mine density \\
        \bottomrule
      \end{tabular}
    \end{table}
    \item \textbf{Safe first move}: if the first action hits a mine, the
      board is regenerated from the same RNG (not reseeded) until the
      first move is safe. This avoids penalizing the agent only due to \textit{its
      bad luck}.
    \item \textbf{Mines never revealed to the agent}: even at a losing
      terminal step, the mine stays $-1$ only in the internal board; the
      \texttt{player\_board} keeps it as $-2$. Loss is communicated
      through the terminal reward alone, keeping the observation
      alphabet consistent with the planned one-hot
      encoding.
  \end{itemize}
\end{frame}

\begin{frame}{Architecture 1: Fully-convolutional CNN}
  \begin{itemize}
    \item inserire disegno esteso con tutti i blocchi e canali
  \end{itemize}
  \begin{highlightbox}{Why fully convolutional?}
    A $1{\times}1$ conv head produces one scalar per cell $\Rightarrow$ the
    same weights work for any $H \times W$ board. A feed-forward classification
    head would have the issue to create space for a $W$ matrix whose shape
    depend on the board size.
  \end{highlightbox}
\end{frame}